% ===================================================================
%  BẢNG Số 11 — Thành phố và các công trình. --- Đám cháy.
%
%  Traduction, faite en 2026, en vietnamien standard d'aujourd'hui,
%  sur l'ido de texto/io. Regles de la colonne : voir 10-bang-01.tex.
%  Deux parties, deux numerotations : les cles portent c1 et c2.
%
%  LA VILLE EST SINO-VIETNAMIENNE DE PART EN PART, ET C'EST ATTENDU :
%  thành phố la ville, quận le quartier, đại lộ le boulevard, công
%  viên le parc, thư viện la bibliotheque, bảo tàng le musee, bệnh
%  viện l'hopital, nhà hát le theatre, tòa thị chính l'hotel de ville.
%  Pas un de ces mots n'est ancien dans la langue parlee : ce sont
%  des composes chinois, et ils sont entres avec la chose qu'ils
%  nomment, entre 1880 et 1930.
%
%  MAIS LE FEU, LUI, EST VIETNAMIEN — et c'est la ligne du tableau.
%  Lửa le feu, khói la fumee, cháy bruler, thang la echelle, vòi le
%  tuyau : la seconde partie, celle de l'incendie du theatre, se lit
%  avec des mots que tout le monde avait avant l'ecole. Il n'y a que
%  les CORPS CONSTITUES qui restent chinois — lính cứu hỏa les
%  pompiers, cảnh sát la police, hiến binh les gendarmes — parce que
%  ce sont des institutions, non des choses.
%
%  QUATRE INSTITUTIONS FRANCAISES RESTENT FRANCAISES, COMME AU
%  QUEBEC : la prefecture, le prefet, les gendarmes et l'ecole
%  normale. Le Vietnam de 1926 les avait, lui — c'etait le meme
%  Etat — et il les nommait « tòa sứ », « quan sứ », « lính khố
%  xanh », « trường sư phạm ». On garde ces noms-la : ce sont ceux
%  que la chose portait ici.
%
%  ET LE WATTMAN DU QUEBEC, QUI Y EST DEVENU CONDUCTEUR DE TRAMWAY,
%  est ici « người lái xe điện », celui qui conduit la voiture
%  electrique. Le mot n'a jamais existe en vietnamien : la chose est
%  arrivee sans lui.
% ===================================================================

%%K t11-apar-1 apar
\VUsaut{2.0mm}
\VUcentre{13.2pt}{90}{LOẠT THỨ BA}
\VUsaut{3.0mm}
\VUfilet{16mm}
\VUsaut{5.6mm}
\VUtitre{15.0pt}{60}{BẢNG Số 11}
\VUsaut{1.4mm}
\VUpk{11.4pt}{40}{Thành phố và các công trình. --- Đám cháy.}
\VUsaut{2.2mm}

%%K t11-tit-1 sub
\VUpk{10.2pt}{40}{Khu ngoại ô.}
\VUsaut{3.0mm}

%%K t11-c1-01-1 p
1. --- Em ở một thành phố lớn cách Đại dương 100 ki-lô-mét, thế mà nó vẫn là một
\VUgras{cảng} \textsuperscript{(1)} hạng nhất. Con sông chảy ven thành phố rộng 400 mét
và tạo thành một vũng đậu rất đẹp.

%%K t11-c1-02-1 p
2. --- Cả phần thành phố nằm trên các \VUgras{bến} \textsuperscript{(2)} mang hẳn dáng
vẻ của biển và của công nghiệp, và làm thành một \VUgras{khu ngoại ô}
\textsuperscript{(3)} chỉ có thủy thủ và thợ thuyền ở. Những người này làm việc trong
các \VUgras{nhà máy} \textsuperscript{(4)} và ở \VUgras{nhà máy khí đốt}
\textsuperscript{(5)}, mà \VUgras{ống khói} \textsuperscript{(6)} cao và \VUgras{bình
chứa khí} nhìn thấy từ rất xa.

%%K t11-c1-03-1 p
3. --- Thời Trung cổ, thành phố có thành lũy, và người ta còn thấy vài dấu tích tường
thành, nhất là \VUgras{cổng} \textsuperscript{(8)} của \VUgras{tòa thành cổ}
\textsuperscript{(7)}, hai bên là hai \VUgras{tháp} \textsuperscript{(9)} nay dùng làm
\VUgras{nhà tù}.

%%K t11-c1-04-1 p
4. --- Trong cả \VUgras{khu phố} trông rất cũ kỹ ấy, chỉ một công trình là tương đối
mới : đó là \VUgras{sở quan thuế} \textsuperscript{(10)}. Nó nằm giữa bến và
\VUgras{con phố} \textsuperscript{(11)} nhỏ dẫn đến \VUgras{tòa thị chính}
\textsuperscript{(12)} cũ. Người ta nhận ra công trình sau này rất dễ nhờ chiếc
\VUgras{tháp chuông} \textsuperscript{(13)} khổng lồ vươn lên trên khu phố cổ.

%%K t11-c1-05-1 p
5. --- Bên kia phố \VUgras{sừng sững} một tòa nhà xây cùng lối với sở quan thuế : đó là
\VUgras{Sở Giao dịch} \textsuperscript{(14)}. Một mặt của nó nhìn ra một quảng trường nhỏ
có \VUgras{tòa sứ} \textsuperscript{(15)}. Thật ra thành phố chúng em tuy không phải thủ
đô nhưng vẫn là một tỉnh lỵ quan trọng.

%%K t11-tit-2 sub
\VUsaut{5.0mm}
\VUpk{10.2pt}{40}{Khu trên của thành phố.}
\VUsaut{3.4mm}

%%K t11-c1-06-1 p
6. --- Đội đồn trú khá đông được đóng trong những \VUgras{doanh trại}
\textsuperscript{(16)} rộng rãi và rất thoáng; chúng trải đến tận hàng rào sắt của
\VUgras{công viên thành phố} \textsuperscript{(17)}. \VUgras{Đại lộ}
\textsuperscript{(19)} dẫn đến công viên ấy chạy sau \VUgras{quỹ tiết kiệm}
\textsuperscript{(18)} và ra tới quảng trường \VUgras{nhà thờ lớn}
\textsuperscript{(20)}.

%%K t11-c1-07-1 p
7. --- Tất cả những công trình ấy xếp thành bậc như một giảng đường trên một quả đồi, nơi
cũng có \VUgras{trường trung học} \textsuperscript{(21)} rộng lớn của chúng em.

%%K t11-c1-07-2 p
Nếu đi xuống theo một con phố hơi dốc, nối vào Phố Lớn, ta đến \VUgras{tòa án}
\textsuperscript{(22)}, trước tòa án trải ra một \VUgras{vườn hoa công cộng}
\textsuperscript{(26)} dễ chịu. \VUgras{Vườn hoa} ấy không lớn lắm, nhưng nó làm thành
giữa lòng thành phố một ốc đảo xanh mát. Vì thế dân thành phố rất hay lui tới, họ thích
đến ngồi dưới mái che của \VUgras{quán cà phê} \textsuperscript{(25)} để nghe các nhạc
công nhà binh chơi vào thứ năm và chủ nhật trong \VUgras{nhà kèn}
\textsuperscript{(27)} đặt giữa các bãi cỏ và bồn hoa.

%%K t11-c1-08-1 p
8. --- Trên một bãi cỏ ấy dựng một \VUgras{bức tượng} \textsuperscript{(28)} bằng đồng,
đài kỷ niệm dựng cho một người đồng hương đã có công lớn với quê hương mình.

%%K t11-tit-3 sub
\VUsaut{4.0mm}
\VUpk{10.2pt}{40}{Phương tiện đi lại.}
\VUsaut{3.4mm}

%%K t11-c1-09-1 p
9. --- Thành phố chúng em chưa được thắp sáng bằng điện, nó mới chỉ có \VUgras{đèn khí}
\textsuperscript{(29)}. Nhưng \VUgras{báo chí ở đây} đang vận động để cải thiện cách thắp
sáng ấy, \VUgras{cũng như} người ta đã cải tiến \VUgras{cách kéo} \VUgras{tàu} điện.
\VUgras{Những} toa xe cũ \VUgras{do} \VUgras{ngựa} kéo đã gần như được thay bằng những
\VUgras{tàu điện tiện lợi} \textsuperscript{(31)}, chở khách nhanh chóng từ đầu này đến
đầu kia \VUgras{của thành phố}, với giá rất \VUgras{rẻ}.

%%K t11-c1-10-1 p
10. --- Gần tòa án là \VUgras{ngân hàng} \textsuperscript{(32)}, nơi người ta chiết khấu
thương phiếu. \VUgras{Người thu tiền} \textsuperscript{(33)} đi thu nợ tận nhà.

%%K t11-tit-4 sub
\VUsaut{4.0mm}
\VUpk{10.2pt}{40}{Nghệ thuật và học vấn.}
\VUsaut{3.4mm}

%%K t11-c1-11-1 p
11. --- Tòa nhà đẹp đứng gần đó là \VUgras{viện bảo tàng}. Nó chứa cùng lúc \VUgras{thư
viện} \textsuperscript{(34)} của thành phố, những \VUgras{bộ sưu tập khoa học tự nhiên}
\textsuperscript{(35)} đẹp đẽ và một \VUgras{phòng tranh} \textsuperscript{(36)} xinh xắn
làm thành một \VUgras{cuộc trưng bày} thường trực tuyệt vời.

%%K t11-c1-11-2 p
Nó mở cửa cho công chúng hai lần một tuần, nhưng khách phương xa có thể vào bất cứ ngày
nào bằng cách cho \VUgras{người gác} \textsuperscript{(37)} một khoản tiền. Các
\VUgras{họa sĩ} \textsuperscript{(38)} đến đó làm việc. Ta thấy họ \VUgras{trước} giá vẽ
của mình, tay cầm \VUgras{bảng màu}, sao lại những bức tranh nổi tiếng.

%%K t11-c1-12-1 p
12. --- Trên cao, sau viện bảo tàng, ta thấy \VUgras{mái vòm} \textsuperscript{(40)} của
\VUgras{bệnh viện} \textsuperscript{(39)} và những tòa nhà của \VUgras{trường sư phạm}
\textsuperscript{(41)}, nơi đã đào tạo các thầy cô của trường làng chúng em. Trên quảng
trường Nhà hát có một \VUgras{bưu cục} \textsuperscript{(43)} đặt trong một \VUgras{ngôi
nhà tư} \textsuperscript{(42)}.

%%K t11-tit-5 sub
\VUsaut{5.0mm}
\VUpk{11.4pt}{40}{Đám cháy.}
\VUsaut{3.0mm}

%%K t11-tit-6 sub
\VUpk{10.2pt}{40}{Cháy!}
\VUsaut{3.4mm}

%%K t11-c2-01-1 p
1. --- Một tin đồn khủng khiếp lan khắp thành phố : nhà hát vừa bốc cháy! Lúc em đến
\VUgras{quảng trường} \textsuperscript{(96)}, đám đông đã chen chúc trên \VUgras{vỉa hè}
\textsuperscript{(46)} và trên \VUgras{lòng đường} \textsuperscript{(47)}. Các
\VUgras{nhân viên bưu điện} \textsuperscript{(44)} và \VUgras{người đưa thư}
\textsuperscript{(45)} chạy ra khỏi bưu cục. Một \VUgras{người lái xe điện}
\textsuperscript{(48)} phải dừng xe để nhường đường cho \VUgras{xe cứu thương}
\textsuperscript{(50)} của thành phố, chở đến một \VUgras{bác sĩ phẫu thuật}
\textsuperscript{(52)} và một \VUgras{y tá} \textsuperscript{(51)} để chăm sóc một
\VUgras{người bị thương} \textsuperscript{(53)} vừa được khiêng trên \VUgras{cáng}
\textsuperscript{(54)} đến gần \VUgras{ki-ốt bán báo} \textsuperscript{(55)}.

%%K t11-c2-02-1 p
2. --- Một đại đội bộ binh đang trên đường tập về đã dừng lại để giữ trật tự cùng
\VUgras{đội tuần cảnh cưỡi ngựa} \textsuperscript{(61)}. Hơn cả \VUgras{lính}
\textsuperscript{(57)} hay \VUgras{hiến binh} \textsuperscript{(63)},
\VUgras{những kỵ sĩ} có những con ngựa chồm lên ấy đẩy lùi được đám \VUgras{người hiếu
kỳ} \textsuperscript{(56)}. Vả lại hiến binh đang rất bận. Một người trong họ chỉ cho
\VUgras{cảnh sát} \textsuperscript{(64)} chỗ phải đưa một \VUgras{nạn nhân}
\textsuperscript{(65)} khác đến, một người lính cứu hỏa dũng cảm ngã từ thang xuống.

%%K t11-c2-03-1 p
3. --- \VUgras{Viên sĩ quan} \textsuperscript{(66)} chỉ rõ chỗ phải đặt chiếc \VUgras{máy
bơm hơi nước} \textsuperscript{(67)} đang tới. Trong khi chờ cỗ máy mạnh mẽ ấy, ông đã
cho lắp một \VUgras{máy bơm tay} \textsuperscript{(68)}, mà chiếc \VUgras{vòi}
\textsuperscript{(69)} dài phun nước như thác vào lửa. Sáu người lính cứu hỏa vận hành
nó, trong khi người bạn của họ cầm \VUgras{đầu vòi} \textsuperscript{(70)} hướng
\VUgras{tia nước} \textsuperscript{(71)} vào giữa đám cháy.

%%K t11-c2-04-1 p
4. --- Phía bên kia nhà hát, người ta đã dựng chiếc \VUgras{thang}
\textsuperscript{(72)} lớn. Nó cho phép lính cứu hỏa đến gần những ngọn lửa đang phụt lên
giữa những cuộn \VUgras{khói} \textsuperscript{(75)} đen.

%%K t11-tit-7 sub
\VUsaut{5.0mm}
\VUpk{10.2pt}{40}{Những hành động dũng cảm.}
\VUsaut{3.4mm}

%%K t11-c2-05-1 p
5. --- \VUgras{Đám cháy} \textsuperscript{(74)} bùng lên ở phòng chờ của \VUgras{nhà hát}
\textsuperscript{(73)}, trong lúc người ta đang tập vở \VUgras{nhạc kịch} mà
\VUgras{buổi diễn} đầu tiên sẽ là ngày mai. May thay chỉ có ít \VUgras{khán giả}
\textsuperscript{(76)} trong phòng. Những người khốn khổ ấy đã lánh ra \VUgras{ban công}
\textsuperscript{(77)}, nơi những \VUgras{lính cứu hỏa} \textsuperscript{(86)} dũng cảm
đến cứu họ bằng thang và dây chão.

%%K t11-c2-06-1 p
6. --- Dưới \VUgras{hàng cột} \textsuperscript{(78)}, nơi tấm \VUgras{áp phích}
\textsuperscript{(79)} báo buổi diễn, ta thấy các \VUgras{nhạc công}
\textsuperscript{(81)} của \VUgras{dàn nhạc} \textsuperscript{(80)} chạy ra, mang theo
nhạc cụ của mình : \VUgras{vĩ cầm} \textsuperscript{(81)}, \VUgras{sáo}
\textsuperscript{(82)}, \VUgras{trung hồ cầm} \textsuperscript{(83)}, \VUgras{kèn
trombon} \textsuperscript{(84)}, \VUgras{trống} \textsuperscript{(85)}, v.v. Người ta cố
cứu một phần \VUgras{đồ đạc} \textsuperscript{(87)}.

%%K t11-c2-07-1 p
7. --- \VUgras{Nam diễn viên} \textsuperscript{(89)} và \VUgras{nữ diễn viên}
\textsuperscript{(88)}, \VUgras{ca sĩ nam} và \VUgras{ca sĩ nữ} đã phải chạy ra với
\VUgras{trang phục} \textsuperscript{(90)} sân khấu của mình. Người hát giọng nam cao kể
cho một \VUgras{nhà báo} \textsuperscript{(92)} nghe sự việc đã xảy ra thế nào.
\VUgras{Phóng viên} đã chạy đến ngay từ phút đầu cùng các nhà chức trách : \VUgras{quan
sứ} \textsuperscript{(91)}, \VUgras{thị trưởng} \textsuperscript{(93)}, \VUgras{cảnh sát
trưởng} \textsuperscript{(94)} và một \VUgras{quan tòa} \textsuperscript{(95)}, những
người sẽ điều tra để quy trách nhiệm.
\VUsaut{6.0mm}
\VUfilet{22mm}
